Out of the three classic areas that focus on automatically constructing control rules from a specification,  Discrete Event Control \cite{ramadge_supervisory_1987}, Reactive Synthesis \cite{pnueli_synthesis} and Automated Planning \cite{nau2004automated}, this paper is more related to the first two.






Discrete Event Control was first formulated for *-regular languages \cite{ramadge_supervisory_1987} and then extended to $\omega$-regular languages. In \cite{thistle_control_1995}, Thistle addresses control problems described using a deterministic finite automaton (DFA) equipped with two acceptance conditions for $\omega$-languages, one representing a specification and the other representing an assumption. Supervisors 
are required to generate infinite words that belong to the specification $\omega$-language under the suposition that the environment generates finite traces of the DFA and infinite traces of the 
assumption $\omega$-language. The supervisor is also required to not block the environment from generating some trace in the assumption $\omega$-language.
In \cite{majumdar_supervisory_2021}, this is referred to as a \textit{non-conflicting} requirement. 
%that every valid finite word generated by the supervisor can be extended to an infinite word in the $\omega$-regular plant assumptions. 
In \cite{thistle_control_1995}, Thistle presents a supervisor synthesis algorithm that is polynomial in the number of states of the automaton and exponential in the acceptance condition size.%for $\omega$-regular specifications under mentioned general $\omega$-regular plant assumptions. 

The control problem we study in this paper can be presented as an instance of the general supervisory control problem for $\omega$-regular languages where the $\omega$-regular language for the  specification is given as an LTL formula. Whereas the $*$-regular language for assumptions is given by a composition of LTS (i.e., the plant) and the $\omega$-regular language for assumptions is the entire set of infinite traces of the plant. 


%Some confusion may arise with the term assumption. In GR(1) control problems, the antecedent of the GR(1) formula is referred to as assumption. It is possible for a controller to ensure the formula by forcing the assumption not to hold. This is commonly considered undesirable. Methodologically, it is appropriate to establish if the negation of the assumptions cannot be enforced by a controller \cite{DBLP:journals/tosem/DIppolitoBPU13}. If this is the case, the assumptions should be changed. A variant of GR(1) where the controller is required to avoid blocking the assumptions even if it can has been proposed \cite{majumdar_environmentally-friendly_2019}.





%In this paper, we define safety modelling assumptions implicitly using automata transition rules and liveness assumptions are defined as the limit of assumptions safety language. Therefore, liveness assumptions becomes true trivially. On the other hand, accepted language condition will be defined as a GR(1) formula. This acceptance condition lead us to a polinomially algorithm to get a supervisor ... \hg{justificar porqué gr1?}

Our work is related to reactive synthesis where the control problem is solely expressed by means of a logical formulae, LTL is commonly used. The system-to-be is a reactive module that shares state variables with the environment. Assumed environment behaviour is also described by means of logical formulae. The underlying interaction model between the controller and the environment differs from that of supervisory control. Reactive synthesis assumes a turn based interaction where each part can change in their turn all the variables it controls. In supervisory control, there are no turns and the environment may win all races. 


The relation between supervisory control and reactive synthesis has received much attention in the last years \cite{ehlers_supervisory_nodate}. Notably, \cite{majumdar_supervisory_2021} shows that the non-conflicting/ non-blocking  requirement of supervisory control requires solving games, obliging games, of higher complexity than those classically used for reactive synthesis. The control problem we study does not suffer from this as the non-conflicting requirement is trivially satisfied. 




The class of GR(1) LTL formulae have been studied extensively as controllers for them can be computed polynomially~\cite{Piterman:2006:SRD} as opposed to the double exponential complexity of general LTL~\cite{pnueli_synthesis}. GR(1) has been studied extensively as a specification language for reactive systems (e.g.~\cite{menghi_specification_2019, DBLP:journals/tse/MenghiTAPVCG23}) and various variants have been proposed (e.g.,~\cite{majumdar_environmentally-friendly_2019}). 

 GR(1) control of DES as defined in this paper was originally proposed in~\cite{DIppolito2008}, a monolithic synthesis algorithm is defined. 
Later, on-the-fly composition construction has been proposed~\cite{Ciolek2020,DBLP:conf/aips/DelgadoSBU23}.  

The reactive synthesis community has explored compositionality as a way to scale synthesis by identifying and performing independent synthesis tasks~\cite{DBLP:conf/atva/FiliotJR10,DBLP:journals/isse/FinkbeinerGP23,DBLP:conf/atva/FinkbeinerP20}. 
%The approach proposed in~\cite{DBLP:conf/cav/KupfermanPV06}  reuses parts of the already existing solution but does not perform independent synthesis tasks. 
%Such Independence in the reactive synthesis literature leverage exploiting . 
For instance, in~\cite{DBLP:conf/atva/FiliotJR10} LTL synthesis is solved compositionally by first computing winning strategies for each conjunct that appears in the large formula using an incomplete heuristic. The decomposition algorithm of~\cite{DBLP:journals/isse/FinkbeinerGP23} works for LTL and decomposes the formulae into  subspecifications that do not share output variables. A dependency graph used for finding the decomposition-critical propositions that, in turn, are key to soundness and completeness of the approach. 
%that can be solved independently and that their results are non-contradictory, i.e., that they can be combined into an implementation satisfying the initial specification. 
In~\cite{DBLP:conf/atva/FinkbeinerP20}  decomposition means partitioning system's input and output variables and then applying incremental synthesis of dominant strategies. Dependency graphs must be constructed to identify components and synthesis order that would ensure completeness. On the other hand, our technique works compositionally in terms of environment's safety assumptions (involving input, output and synchronization events) which are described as synchronizing LTS. Goal formula decomposition is a consequence of projection of the formula onto intermediate subsystems. Grouping and ordering of compositions does not affect neither soundness nor completeness (although it may affect performance).

%A more fine grained approach taking into account dependencies is presented in~\cite{DBLP:journals/isse/FinkbeinerP20} but authors admit greater overhead  than their later approach~\cite{DBLP:journals/isse/FinkbeinerGP23}.


Traditional compositional verification~\cite{DBLP:conf/cav/GrafS90,DBLP:conf/popl/ClarkeGL92} mitigates the state-space explosion by leveraging modular structure of system-under-verification. In that vein, compositional synthesis for plants expressed as finite-state systems has been studied by the community of supervisory control for non-blocking goals~\cite{DBLP:journals/deds/FlordalMFA07,DBLP:journals/deds/MalikMF23}. There, synthesis is  addressed by incrementally composing (and reducing) plant's component. Also, resulting controllers are modular: they are  made up by automata that interact by mean of shared labels.  Our solution is inspired in such works but it addresses  intfinite trace goals for which there is no guarantee for the existence of maximal solutions.

%Related to compositional approaches and modular solutions is the problem of distributed synthesis~\cite{DBLP:conf/focs/PnueliR90}: yielding decentralized controllers that are supposed to run on different agents that have a partial observation of the environment. The general problem is known to be undecidable~\cite{DBLP:conf/focs/PnueliR90}. Thus, most works, in either discrete or hybrid controllers, either fix how the finite-state environment model looks like (e.g., grid with moving obstacles, discretizations of continuous dynamical system, etc.) (e.g.~\cite{DBLP:journals/iandc/AlurMT18}),~\cite{DBLP:journals/tcad/MajumdarMSZ20,DBLP:journals/iandc/AlurMT18,DBLP:journals/arobots/MoarrefK20}, work for certain types of local control specifications~\cite{DBLP:conf/icalp/MadhusudanT01,DBLP:journals/tac/MallikSSM19,DBLP:journals/tac/PolaPB18,DBLP:conf/cdc/BoskosD15}, safety specifications~\cite{DBLP:journals/tac/MeyerGW18} or attractivity~\cite{DBLP:conf/cdc/Apaza-PerezG21} and, fundamentally, either require an assume-guarantee/abstraction (or small-gain)  framework, bounds or locality assumptions at expense of general completeness~\cite{DBLP:journals/arobots/MoarrefK20,DBLP:conf/cdc/DallalT15,DBLP:conf/cdc/Apaza-PerezG21,DBLP:journals/tac/PolaPB18,DBLP:conf/cdc/BoskosD15,DBLP:journals/isse/FinkbeinerP22}, or redefine the synthesis problem as a non-zero sum game~\cite{DBLP:conf/tacas/ChatterjeeH07}.
%In such approaches, a system with a control specification is decomposed into subsystems with local control specifications. Then, for each subsystem, a symbolic abstraction can be computed and a local controller is synthesized while assuming that the other subsystems meet their local specifications.[28], [24](reisig), [19] P.-J. Meyer, A. Girard, and E. Witrant. Safety control with performance guarantees of cooperative systems using compositional abstractions. 
%Although we produce modular controllers, we do not aim at decentralization and thus we do not  require neither to sacrifice completeness.



%\footnote{Although GR(1) strategies are not truly memoryless, memory only serves as a way to systematically traverse goals. Changing a finite number of times the memory associated to a plant state would not affect strategy reaching recurrently goals \su{Creo que el problema es este footnote, no está el background para entenderlo. Creo que debería ir en related work o discussion, %\vb{Acá est'a la clave del asunto. Lo que tiene que pasar es tener una forma efectiva pero fundamentalmente sound para saber que un determinado estado de la planta NO puede ser parte de ninguna estrategia ganadora DESDE EL ESTADO INICIAL. Se puede asumir que uno tiene un predicado en la regla en cuestión y habría explicar qué hacemos, por ejemplo, para GR[1] en particular (es la parte de safety la que usamos?).  }\su{Sí, esto debería decirse más abajo. Donde?}\su{Reformulé el texto para no hablar de estrategias y dar pistas sobre como se va a usar winning states}








